% ================================================================
%  第 6 章 梦与做梦者 (原书 p.133-146, 前半: 6.1-6.8)
%  源文: _work/src1/task-ch06a.md
%  说明: 本章后半 (p.147-152, 6.9-6.11) 见 part1-ch06b.tex
% ================================================================
\chapter{梦与做梦者}

\noindent 1984年2月\par

% ---------------- p.133 ----------------
\section{(5) 被禁的梦}

\origpage{9} 借着《追寻栈》(À la Poursuite des Champs)写作中断三个月的机会,我要把《导论》接着去年六月停下的地方写下去。相隔半年多,我刚把它仔细重读了一遍,并添了几个小标题。

写这篇《导论》时,我很清楚:这类思考免不了招来一大堆“误解”——而想事先堵住它们也是白费力气,那无非是在原有的误解之上再堆上一些!关于这一点,我只想补充一句:我丝毫无意向那种由千年惯例确立起来的科学写作风格开战——我自己就曾勤勤恳恳地奉行它二十多年,还把它当作数学家这一行当的本质部分教给我的学生们。不管对错,今天我也仍然这样看待它,并且继续这样教。甚至可以说,我倒更显得老派:我坚持工作要做到底,从头到尾一针一线地缝,不放过任何一个稍有阴影的角落。近十年来我不得不往酒里掺水,那实在是形势所迫!“规范的写作”在我看来仍然是数学工作的重要一环,既是一种发现的工具——用来检验并深化对事物的理解,没有它,这理解始终是粗略而零碎的——也是一种传达这种理解的手段。从教学的角度看,严格的表述方式,也就是演绎的方式(它绝不排除铺开宏阔画卷的可能),有明显的好处:简洁,查考方便。这些好处是实实在在、有分量的——当讲稿面向
% ---------------- p.134 ----------------
\origpage{10} 数学家(比如说)、而且是已经对所述主题的某些来龙去脉、或对与之十分相近的主题足够熟悉的数学家时。

然而,若一篇讲稿的对象是孩子、是年轻人,或是那些事先完全不“在行”的成年人——他们的兴趣已经觉醒,而且多半(并将继续如此,而且事出有因……)对一项发现性工作的真正进程全然无知——那么这些好处就完全成了幻影。更确切地说,是这样的读者:他们连这种工作的存在都不知道——而这工作其实人人做得到,只要他有好奇心和常识;我们对宇宙万物的智性认识,正是从这工作中不断生出、又不断再生,包括那些以宏伟的编排体系表达出来的认识,比如欧几里得的《几何原本》,或达尔文的《物种起源》。对这种工作的存在与性质全然无知,几乎是普遍现象,各级教师也不例外,从小学教师到大学教授。这是一个非同寻常的事实;去年写这篇《导论》的第一部分、开始这番反思时,它才清清楚楚地显现在我眼前——同时也让我隐约看到了这件令人困惑之事深深的根……

即便讲稿面对的读者在各方面都确实“在行”,仍有一件重要的东西,是“严格的”表述方式不许自己传达的。而这东西,在严肃人士的圈子里——尤其在我们这些科学家当中——也完全不受待见!我说的就是梦。梦,以及梦向我们耳语的那些景象——起初像梦一样不可触摸,而且往往迟迟不肯成形。长年累月、甚至整整一生的紧张工作,也许都不足以让某个梦的景象完全显现出来,不足以看它凝聚、打磨,直到获得钻石的硬度与光华。这就是我们的工作,无论我们是用手还是用精神做工的工匠。当工作完成、或工作的某一部分完成时,我们把那有形的成果呈现在所能找到的最亮的光下,我们为它欢喜,还常常以它为荣。然而,我们长久打磨的这颗钻石里,并没有那个在打磨中激励我们的东西。也许我们造出了一件极精确的工具、一件好用的工具——但工具本身就是有限的,就像人手所做的一切,哪怕它在我们看来很了不起。一个景象,起初没有名字、没有轮廓,淡薄得像一片雾霭,它引导着我们的手,让我们伏在活儿上,感觉不到时辰流逝,也许也感觉不到年月。一片从雾与幽暗的无底之海中无声地剥落下来的雾霭……我们身上那无界限的东西,就是她,就是这片海——只要我们的渴望使她受孕,她便随时准备孕育、随时准备分娩。梦正是从这场婚配中涌出,像蜷伏在养育它的子宫里的胚胎,等待着幽暗的劳作把它带向第二次诞生,带向天光。

一个蔑视梦的世界有祸了——那也是我们身上深刻的东西遭蔑视的世界。我不知道在我们这种文化(电视、电脑和洲际火箭的文化)之前,是否还有别的文化宣扬过这种蔑视。这大概正是我们区别于前人的许多地方之一——我们把前人彻底取代、几乎可以说从地球表面清除掉了。我不曾知道有哪一种别的文化是不尊重梦的,不曾知道有哪一种文化里,梦的深根不被所有人感受到、不被承认。在一个人或一个民族的一生中,可曾有哪件宏大的作品,不是生于梦、不是先由梦哺育,而后才在大白天里绽放的?然而在我们这里(是不是已经该说:到处都一样?),对梦的敬重被叫作“迷信”;而众所周知,我们的心理学家和精神病学家早已把梦翻来覆去、量了个遍——那点分量,想必勉强够塞满一台小型电脑的存储器。也确实,“在我们这里”,再没有人会生火,也没有人敢在自己家里看着孩子出生,或看着母亲、父亲死去——
% ---------------- p.135 ----------------
\origpage{11}有诊所和医院就是干这个的。谢天谢地……我们这个世界,如此自豪于它的力量——以原子百万吨级计,以图书馆和电脑里储存的信息量计——恐怕也正是这样一个世界:在这里,每个人的无能,面对生活中那些简单而本质之物时的恐惧与蔑视,达到了顶点。

\origpage{12} 幸好,梦就像最压抑的社会里那原初的性冲动一样,命硬得很!不管是不是迷信,它仍旧偷偷地、执拗地向我们耳语一种认识——那是我们醒着的精神太笨重、或太怯懦,把握不住的认识——并且给它启发我们形成的种种计划注入生命、插上翅膀。

我刚才似乎暗示过,梦往往不肯成形;那只是表面现象,并未真正触及事情的底。那“不肯”多半来自我们醒着的精神,来自它平常的“心态”——而且“不肯”这个词还是轻描淡写!那更像是一种深深的戒心,它掩盖着一种古老的恐惧——认识的恐惧。就本来意义上的梦而言,梦的消息越是贴近我们,越是因为一旦被听见就会带来我们人格深处蜕变的威胁而显得沉重,这恐惧就越是起作用,它立起的屏障就越是有效。但看来,即便在数学的“梦”这种相对无害的场合,这份戒心也依然在场、依然有效;以至于一切梦似乎不仅被逐出了文本(反正我连一部记下梦的文本也没见过),而且也被逐出了同事间的谈话——不论是小圈子里,甚至一对一的时候。

事情若果然如此,那当然不是因为数学的梦不存在、或不再存在了——若真如此,我们的科学早就贫瘠了,而绝非如此;这种表面上的缺席、这种沉默的合谋,其原因想必与另一种共识紧密相连——那就是要小心抹去一切痕迹、一切提及:抹去那借以完成发现、借以更新我们对世界之认识的工作。或者不如说,围绕着梦、也围绕着梦所唤起、所启发、所滋养的工作的,是同一个沉默。以至于“数学的梦”这个说法本身,在许多人看来都会是无意义的胡话——我们如此经常地被一些按电钮式的陈词滥调驱动,而不是被我们对某种极其简单、日常而又重要的现实的直接经验驱动。

\section{(6) 做梦者}

事实上,经验告诉我:当精神渴望认识梦,而不是逃避它(或手里拿着一副现成的框框去套它——那是一样的),梦就丝毫不会“不肯成形”——它肯让人细致地描述,肯交出它的消息:那消息永远简单,从不愚蠢,有时还叫人震动。恰恰相反,我们心中的“做梦者”是一位无与伦比的大师:他每一次都能找到、或凭空造出最合用的语言,来绕过我们的恐惧、摇醒我们的麻木;他的舞台手段变化无穷,从完全没有一丝视觉或感官要素,到最令人瞠目的场面调度。当他显现时,绝不是为了躲开,而是为了鼓励我们(几乎总是白费力气,而他的好意从不疲倦……)走出我们自己,走出他看见我们深陷其中的那种沉重——有时他还会不动声色地拿这种沉重寻开心,用滑稽的色调把它戏仿一番。倾听我们心中的“做梦者”,就是与我们自己交流——逆着那些强大的堤坝,它们非要禁止我们这样做不可。

但能做大的人,也能做小。既然我们能借梦与自己交流、向自己显现自己,那么,把数学的梦——比如说——那绝不私密的消息传达给别人,想必也同样简单易行;这种梦并不牵动
% ---------------- p.136 ----------------
\origpage{13} 力量相当的抵抗。

而说实话,在我做数学家的过去里,除了追随一片片从浓重致密的雾的织物上一一剥落下来的梦,把它们“梦”到底——直到它们最显明、最坚实的显现,无可辩驳的显现——我还做过别的什么?又有多少次,我为自己那份固执而急得跺脚:我总要把每块凝结着我的梦的宝石、或半宝石,一丝不苟地打磨到最后一个刻面——而不是听从一股更深的冲动:去追随那母体织物千变万化的幽径,直到梦与它显在的、可以“发表”的化身(总之,按通行的规范)之间那片朦胧的边界!况且,我当时正要听从那股冲动,投入一项“数学科幻”的工作——一部关于“Motif”理论的“某种白日梦”,那理论当时还纯属假设,而且直到今天仍然如此,而且理所当然:因为再没有另一个“白日做梦者”投身这场冒险。那是六十年代末;那时我的生活(我丝毫没有想到)正要拐一个大弯,此后十来年间,这转折把我的数学热情挤到了边缘的位置,甚至遭到否认。

但细究起来,《追寻栈》——沉默十四年之后的第一部出版物——恰恰属于那场从未写出的“白日梦”的精神,而它似乎成了那场梦临时的后续。当然,至少初看之下,这两个梦的主题之间的差异,是两个数学主题之间所能有的最大差异;何况头一个(即 Motif 那个)似乎位于“用现有手段可做之事”的地平线之外,而第二个——即著名的“栈”及其同类——看起来完全伸手可及。这些差异,或许可以称为偶然的、意外的差异,它们也许会比人们预料的早得多就烟消云散\fnoteP{1}{(3)}{我在这里特别想到已故的莫德尔(Mordell)猜想、泰特猜想、沙法列维奇(Chafarévitch)猜想——去年,法尔廷斯(Faltings)一篇四十页的手稿把这三个猜想一并证明了;而就在此前,“在行”的人们早已确立的共识是:这些猜想“够不着”!恰好,我那心爱的“远阿贝尔代数几何”纲领当作拱心石的那条基本猜想,正与莫德尔猜想相近。(甚至据说后者可以由前者推出——这就更说明,这个纲领不是给严肃人士看的故事……)} (3)。在我看来,这些差异对于两个主题各自可能引发的这类工作,影响相对很小——因为这里说的正是“白日梦”,或者用不那么挑衅的说法:把概念上的粗加工一直做到一个足够融贯、足够精确的整体景象,足以或多或少完整地使人确信:这景象在根本上确实符合事物的实际。就本书所展开的这个主题而言,这大致意味着:对这一景象之有效性的逐项验证,变成了一件纯属本行手艺的事。这当然可能要求巨大的工作量,其中也需要机巧和想象,而且想必还有意外的转折与前景,幸而会使它不至于成为纯粹的例行公事(用安德烈·韦伊的话说,一场“长长的练习”)。

总之,这正是我过去做腻了又重做的那类工作,它就在我指尖上,因此在我余下的岁月里不必再做。只要我再度投入数学工作,我的精力肯定最该用在“白日梦”的边缘地带。在这个选择里,启发我的并不是什么效益考虑(假如这种考虑能启发得了谁的话),而恰恰是一个梦,或一些梦。我身上这股新的冲劲若能显出力量,那力量必定是从梦里汲取的!

% ---------------- p.137 ----------------
\section{(7) 伽罗瓦的遗产}

\origpage{14} 看起来,在所有的自然科学中,只有数学把我所谓的“梦”、或“白日梦”判以一道看来绝对的禁令,而且这禁令已施行两千多年。在别的科学里,包括物理学这样公认“精确”的科学,梦至少是被容忍的,甚至(视时代而定)是被鼓励的,只不过用的是些更“拿得出手”的名字,比如:“思辨”、“假设”(如著名的“原子假设”,它出自一个梦——对不起,出自德谟克利特的一种思辨)、“理论”……从那个“不敢道出自己名字的梦”的身份,过渡到“科学真理”的身份,是经由几乎察觉不到的层层台阶、经由一个逐渐扩大的共识完成的。相反,在数学里,这几乎总是(至少如今)一次突变,靠的是一篇证明的魔杖一挥\fnoteP{2}{(4)}{况且就在今天,我们也还遇得到身份可疑的“证明”。格劳尔特(Grauert)那个以他命名的有限性定理,多年来就是这样:他那个证明谁也不曾读完(而愿意读的人并不缺!)。这一困窘后来由另一些更透明的证明解决了,其中有些还走得更远,它们取代了最初那个证明。类似而更极端的情形,是所谓的“四色问题”的“解决”:其中计算的部分是用电脑(以及几百万美元)了结的。所以那是一个“证明”,它的依据不再是理解一个数学情境而来的内心确信,而是对一台没有理解能力的机器的信任——而使用它的数学家既不知道这台机器的结构,也不知道它如何运作。即便计算再由别的电脑、按别的计算程序核实一遍,我也不认为四色问题就此了结。它只是换了一副面孔:如今要做的几乎不再是找一个反例,而只是要一个证明(当然,得是读得懂的证明!)。} (4)。在那样一些年代——数学定义与数学证明的概念还不像今天这样清楚,还不成其为(或多或少)普遍的共识——却也有些显然重要的概念,身份暧昧不明,比如“负”数(帕斯卡拒斥过它),或“虚”数。这种暧昧,至今仍反映在通行的语言里。

\origpage{15} 定义、陈述、证明、数学理论这些概念的逐步澄清,在这方面极为有益。它让我们意识到我们手头的工具——尽管简单得像小孩玩意——有多么强大:靠它们,我们能把那看似无法表述的东西表述得完美精确——几乎只要足够严格地使用日常语言就行。如果说有什么东西从我童年起就让我着迷于数学,那恰恰是这种力量:用词语去圈住、并完美地表达某些数学事物的本质——这些事物乍看之下形态如此飘忽、如此神秘,仿佛在言语之外……

然而这种力量、这种完美精确与证明所提供的资源,有一个糟糕的心理上的反作用:它们更加强化了对“数学的梦”的传统禁忌;也就是说,强化了对一切不肯以常规的精确面貌出现的东西的禁忌(哪怕牺牲的是更宏阔的景象)——那种精确要有形式完整的证明来担保它“货真价实”,否则(而且如今越来越如此……)就得有据信可以整理成形的证明概要来担保。偶尔的猜想,勉强可以容忍,条件是它得符合问卷式的精确要求,而问卷上只允许“是”或“否”两种回答。(而且,还用说吗,还得是那个敢提猜想的人在数学界里有门有面。)据我所知,还没有过一个例子:把一门在其本质部分明明白白是猜想性的数学理论,当作“实验”来发展。确实,按现代规范,十七世纪以来发展起来的整套“无穷小”演算——后来成了微分学与积分学——会显得像一场白日梦,而它最终要等到两个世纪之后,才靠柯西的魔杖一挥,变成
% ---------------- p.138 ----------------
严肃的数学\origpage{16}。这必然让我想起埃瓦里斯特·伽罗瓦(Evariste Galois)的白日梦——他不巧也碰上了同一位柯西,运气不佳;但这回不到一百年就够了:又一下魔杖(这回是若尔当,如果我没记错),就让这场梦获得了公民权,并为此改名为“伽罗瓦理论”。

从这一切得出的结论——这结论对“1984年的数学”并不有利——是:幸好牛顿、莱布尼茨、伽罗瓦这样的人(肯定还漏掉很多,我对历史并不内行……)当年没有被我们今天的规范绊住;那时候,他们只管发现,没工夫立规范!

伽罗瓦这个例子不请自来,触动了我的一根敏感的弦。我似乎记得,第一次听人说起他和他的奇崛命运时(那时我大约还在念中学或大学),一种兄弟般的亲切感就在我心里醒来了。和他一样,我感到自己心中有一份对数学的激情——也和他一样,我觉得自己是个边缘人,是那个(在我看来)把他抛弃了的“上流社会”里的异乡人。可我自己最终也成了这上流社会的一员,又在某一天离开它,没有遗憾……这份有点被遗忘的亲近感,最近又以全新的面貌重新出现在我身上——当时我正在写《纲领概要》(Esquisse d'un Programme,为着我申请作为研究员进入国家科学研究中心[CNRS]的事)。这份报告主要概述我近十年来一些主要的思考主题。在所有这些主题中,最让我着迷、也是我打算在今后几年着重发展的那个,正是数学的梦的典型——它还接上了“Motif 之梦”,为后者提供了一条新进路。写这份《草稿》时,我想起了这十四年里我连续做过的、为期最长的一次数学思考。它从1981年1月一直持续到6月,我给它取名叫《穿越伽罗瓦理论的长征》。一针引一线,我意识到:我这几年来时断时续地追着的那场白日梦,后来得了“远阿贝尔代数几何”之名,其实无非是伽罗瓦理论的一个延续,“伽罗瓦理论的终极归宿,而且想必是合乎伽罗瓦本人的精神的”。

当我写到被引的那句话所在的那一段、这种连续性在我眼前出现时,一阵喜悦穿过我,至今未散。那是在完全的孤独中做下来的工作所得到的报偿之一。它的出现,与不久前我从两三位同事兼旧友那里受到的、颇为冷淡的接待同样出乎意料——他们明明是“在行”的人,其中一位还曾是我的学生;我曾有机会跟他们谈起我正在发现的东西,那时我还“热着”,心里满是欢喜……

这让我想到:今天接续伽罗瓦的遗产,必定也要接受他当年那份孤独的风险。也许时代的变化比我们想的要小;不过,这份“风险”在我眼里往往并不成为威胁。我若有时因所爱之人的故作冷淡或轻蔑而难过、失落,那么相反,多年来孤独——数学的或别的孤独——从不曾让我觉得沉重。若有一位忠实的女友,我每次离开她后都不断渴望再回到她身边,那就是她!

\section{(8) 梦与证明}

但让我们回到梦,回到几千年来压在它身上的那道数学禁令。这也许是所有先验之见中最顽固的一条——它往往不明说,扎根在习惯里,裁决着什么东西“算数学”、什么不算。要过几千年,那些像孩童玩意一样简单
% ---------------- p.139 ----------------
而又无所不在的东西——某些几何图形的对称群、另一些图形的拓扑形态、数字零、集合——才获准进入圣殿\origpage{17}!

\origpage{18} 当我跟学生们谈球面的拓扑,谈由球面添上环柄所得的那些形态——这些东西并不叫小孩子吃惊,却叫他们不知所措,因为他们自以为知道“数学”是什么——我得到的第一个自发反应总是:这哪算数学!数学嘛,当然是勾股定理、三角形的高、二次多项式……这些学生并不比你我更蠢;他们的反应,和古往今来、直到今天的世上所有数学家的反应一样,除了毕达哥拉斯或黎曼这样的人,也许还有五六位别的。就连庞加莱——那可不是等闲之辈——也能用一套深有感触的“甲加乙”式哲学,“证明”无穷集合算不上数学!想必也曾有过一个时候,三角形和正方形算不上数学——那只是孩子或陶匠在沙上、在瓶罐的泥坯上画的图,别搞混了……

精神被自己的“知识”憋住,这种根子上的惰性当然不是数学家独有的。我有点偏离正题了:我要谈的是压在数学的梦身上的禁令,以及借由它、压在一切不以制成品的惯常面貌出现(不是开箱即用)的东西身上的禁令。我对别的自然科学所知的那一点,已足以让我估量出:同样严厉的禁令早就会把它们判给贫瘠,或判给乌龟般的进展——有点像中世纪,那时容不得对《圣经》的字句刨根问底。但我也很清楚:发现的深源,以及发现之进程在其一切本质方面,在数学里与在宇宙中我们的身体和精神所能认识的任何别的领域、别的事物里都是一样的。放逐梦,就是放逐源头——判它只能在暗处存在。

我还很清楚(这经验自我少年时代与数学初恋爱慕以来从未被推翻):在展开一种广阔或深刻的数学事物之景象时,正是这景象与理解的展开、这渐进的深入,始终先于证明,使证明成为可能,并赋予它意义。当一个情境——从最卑微到最宏阔——在其本质方面被理解了,对所理解之物的证明(以及对其他东西的证明)就像熟透的果子一样落下来。而像青果一样从认识之树上硬拽下来的证明,会留下一丝不满足的余味,一种渴望未得平息的失落。在我做数学家的一生中,有两三次,我不得不在别无办法时下定决心去拽果子,而不是摘果子。我不是说我做得不好,也不是说我后悔。但我所做得最好的、我所最爱的东西,我都是心甘情愿取来的,不是强夺来的。如果说数学给了我丰盈的快乐,在我成熟之年仍让我着迷,那靠的不是我从它那里硬拽来的证明,而是我在它身上感到的无穷奥秘与完美和谐——它随时准备向一只爱它的手、一道爱它的目光显露自身。

\section{(9) 受欢迎的异乡人}

我觉得时候到了,该谈谈我与数学家世界的关系。这与我与数学的关系完全是两回事。后者从我年少时起就存在,而且很强,远在我想到世上竟有一个数学家的世界、一个数学家的圈子之前。那是一整个复杂的世界,有它的学会、它的期刊、它的聚会、讨论会、大会,有它的头牌女伶和苦力,有它的权力结构,有它的灰衣主教,还有同样灰扑扑的一大群任人使唤、任人摊派的人,为论文或文章发愁;也有另一些更少见的人,他们有的是手段和想法,却
% ---------------- p.140 ----------------
处处碰壁\origpage{19},绝望地找不到一位有权有势者撑腰——那些忙忙碌碌、叫人畏惧的强人握有一种魔法:让一篇文章得以发表……

\origpage{20} 我是在1948年、二十岁上到巴黎时,才发现有一个数学世界的;我那只瘪瘪的手提箱里,装着蒙彼利埃大学的理学学士学位,还有一部行距排得极密、双面书写、不留页边的手稿(纸很贵!),那是我三年孤独思考的记录,思考的东西(后来我才知道)当时早已以“测度论”或“勒贝格积分”之名广为人知。因为从没遇见过别人,我直到进京城那天都还以为,世上只有我一个人在“搞数学”,因此也就是唯一的数学家。(在我眼里这两者是一回事,而且直到今天也还多少是。)我曾玩弄那些我称之为可测的集合(其实也没遇到过一个不可测的集合……),玩弄几乎处处收敛,却不知道拓扑空间是什么。有十来个互不等价的“抽象空间”与紧性的概念,我一直没搞清,那是从一本小册子里捞来的(作者好像叫阿佩尔,收在《科学与工业时事》丛书里),这本小册子是我碰上的——天知道怎么碰上的。那时我还没听人说过——至少没在数学的语境里听过——群、域、环、模、复形、同调(还有别的,就不数了)这类古怪或野蛮的词;忽然之间,它们不打招呼,一齐朝我压下来。这一击不轻!

我之所以在这一击之下“活了下来”,还继续搞数学、甚至以此为业,是因为在那遥远的年头,数学世界还不太像它后来变成的样子。也可能是我运气好,落进了这个未曾料想的世界里一个格外好客的角落。我带着蒙彼利埃学院一位教授的一封含糊的推荐信,那是苏拉先生(跟他的同事们一样,他也没怎么在课堂上见过我!),他当年是嘉当的学生(父亲还是儿子,我已说不太准了)。因为埃利·嘉当(Elie Cartan)当时已经“退出赛场”,他儿子亨利·嘉当(Henri Cartan)便成了我有幸遇见的头一个“同类”。我当时并不知道这是多么吉利的兆头!他接待我时的那份礼貌中透着善意,那是他的标志,几代高师学生都熟悉——他们有幸在他手下初试身手。从他当时给我指点学业方向的那些建议看,他大概没有意识到我无知到了何等地步。无论如何,他的善意显然是冲着人来的,不是冲着行李或可能有的天分,也不是(后来)冲着名声或名气。

接下来那一年,我成了“高师”嘉当一门课上的常客(讲流形上的微分形式体系),我拼命跟;也常去“嘉当讨论班”,目瞪口呆地看着他和塞尔辩论,满嘴“谱序列”(呜!),满黑板的图(叫做“图表”),画满了箭头。那是“层”、“甲壳”以及一整套军火库的英勇时代,那些东西的意义我完全摸不着边,可我还是勉勉强强逼着自己把定义和陈述吞下去,把证明一一验过。在嘉当讨论班上,谢瓦莱、韦伊也定期露面;而到了布尔巴基讨论班的日子(参加者和听众加起来,撑死也就二三十人),就看见这个有名的布尔巴基帮的其余成员——迪厄多内、施瓦茨、戈德芒、德尔萨特——像一伙有点吵闹的哥们儿那样闯进来。他们彼此以“你”相称,说着一种我几乎完全听不懂的共同语言,烟抽得凶,也爱笑;只差几箱啤酒,气氛就齐全了——而啤酒是由粉笔和海绵代替的。那气氛与法兰西公学院勒雷的课截然不同(讲肖德尔关于无限维空间中拓扑度的理论,可怜的我!),我是听嘉当的建议去听的。我曾到法兰西公学院找过勒雷先生,问他(如果我没记错)这门课要讲什么。他可能给我的解释,我一点不记得了,
% ---------------- p.141 ----------------
\origpage{21} 也不记得自己有没有听懂什么——只记得:在那里我也感到一种善意的接待,那是给随便哪个初来乍到的外乡人的。肯定正是这一点,而非别的什么,让我去上那门课,并且像在嘉当讨论班那样硬着头皮跟下去——尽管勒雷在那里讲的东西,我当时几乎完全不懂。

奇怪的是:在这个世界里我是新来的,它的语言我几乎不懂,说得更差,可我并不觉得自己是个外乡人。虽然我几乎没机会(而且当然没本事)跟韦伊、迪厄多内这样快活的家伙说话,或跟嘉当、勒雷、谢瓦莱这样气度更雍容的先生们说话,我却觉得自己被接纳了,几乎可以说:是他们中的一个。我不记得有哪一次被这些人中的某一位居高临下地对待,也不记得有哪一次,我的求知欲——以及后来,又一次,我发现的喜悦——被一份自得或一份轻蔑挡了回去\fnoteP{3}{(5)}{这个事实尤其值得注意:直到1957年前后,布尔巴基小组的成员中不止一位对我还颇有保留,而他们最终把我增选进去时,我相信也带着几分不情愿。一句善意的玩笑话把我归入了“危险专家”(指泛函分析)之列。我有时感到嘉当身上有一种更深的、未说出的保留——有那么几年,我一定给了他这样的印象:一个热衷于空洞浮泛的推广的人。他曾惊讶地发现,我为布尔巴基写的那一篇(也是唯一一篇)稍长的稿子(论流形上的微分形式体系)里,竟有一点相当实在的想法——而我提议由我来写时,他并不怎么热心。(多年以后,在从凝聚对偶的观点发展留数形式主义时,这一想法又派上了用场。)况且在布尔巴基的会议上,我多半是跟不上的,尤其是逐篇共同宣读稿子的时候,我实在没法跟上宣读与讨论的速度。可能我确实不太适合集体工作。不管怎么说,我在融入共同工作上的这份困难,或我因别的缘故在嘉当和别一些人那里可能引起的保留,都不曾在任何时候给我招来讥讽或碰钉子,或哪怕一丝居高临下——至多在韦伊那里有过一两次(此人确实是个特例!)。嘉当对我的态度,从未有一刻离开过一贯的和善,其中带着亲切,也带着他特有的那点幽默——在我,那点幽默始终与他这个人分不开。} (5)。若不是这样,我也不会像人们说的那样“成为数学家”——我会另选一行,在那里我能施展自己,而不必去面对轻蔑……

尽管“客观上”我是这个世界的异乡人,一如我在法国是个异乡人,却有一条纽带把我和这些出身别的阶层、别的文化、别的命运的人连在一起:一份共同的激情。我怀疑,在我发现数学家世界的那个关键年头,他们中没有一个人——连嘉当也没有,我勉强算他的学生,可他还有许许多多别的学生(而且是不那么跟不上的学生!)——在我身上察觉出那份同样盘踞在他们心中的激情。在他们眼里,我大概只是课堂和讨论班一大群听众中的一个,记着笔记,显然不太跟得上。如果说我多少有点与众不同,那就是我不怕提问——而那些问题多半只是暴露出我在语言和数学两方面都无知得惊人。回答可能简短,甚至带着惊讶,但当时那个目瞪口呆的冒失鬼,从未碰过钉子,从未被“放回原位”——无论在布尔巴基那伙人不拘礼数的圈子里,还是在法兰西公学院勒雷课堂上更严峻的氛围里。那些年里,自从我口袋里揣着一封给埃利·嘉当的信来到巴黎,我从未觉得自己面对着一个宗派、一个封闭乃至敌意的世界。如果说我领教过、深谙面对轻蔑时那种内心的紧缩,那不是在那个世界里;至少不是在那时。对人的尊重,是我在那里呼吸的空气的一部分。不必去挣得尊重,不必先证明自己才能被接受、被和气地对待。说来也许奇怪:只要是一个人,只要有一张人的面孔,就够了。

% ---------------- p.142 ----------------
\section{(10) “数学共同体”:虚构与现实}

\origpage{22} 因此毫不奇怪:也许就在那一年,在我内心深处,我已感到自己是这个世界的一员;而在随后的年月里,这种感觉越来越清楚。我喜欢用“数学共同体”这个名称来指它,这名字对我满载着意义。在写下这些字句之前,我从未有机会细想过自己给这名字赋予了什么意思,尽管我在很大程度上认同于这个“共同体”。现在清楚了:它对我而言,不外是那个善意接待我、把我当作他们中一员的世界的理想延伸——在空间上和时间上的延伸;而且是一个与我生命中最主要的激情之一相连的世界。

\origpage{23} 我逐步认同的这个“共同体”,并不是对最初接待我的那个数学圈子所作的完全虚构的外推。最初的圈子一点点扩大了;我的意思是:在最初接触之后的十年、二十年里,由于共同感兴趣的主题和人与人之间的投缘,我经常来往的数学家圈子不断扩大。具体说来,就是同事与朋友的圈子,或者不如说那个同心结构:从与我关系最近的同事(先是迪厄多内、施瓦茨、戈德芒,后来尤其是塞尔,再后来是安德烈奥蒂、朗、泰特、扎里斯基、广中平祐、芒福德、博特、迈克·阿廷这样的人,还不算布尔巴基小组的人——那个小组也在一点点扩大——以及六十年代起投到我这里来的学生们……),直到我在此处彼处有机会遇到、由或强或弱的投缘而或紧或松地联系着的其他同行——总之,正是这个由相遇与投缘偶然凑成的微观世界,构成了那个于我充满温度与回响的名字的具体内容:数学共同体。当我把它认同为一个活生生的、温暖的实体时,实际上我认同的就是这个微观世界。

直到1970年那个“大转折”之后——我该说,那第一次觉醒之后——我才意识到:这个舒适可亲的微观世界只占“数学世界”极小的一部分;而我乐于赋予那个世界、那个我一直不去了解、从没想到要去关心的世界的那些特征,都是虚构的特征。

况且在这二十二年间,这个微观世界本身也换了面孔,而它周围的世界也在变。我自己当然也变了——年复一年,不知不觉,我和我周围的世界一样变了。我不知道我的朋友和同事们是否比我更觉察到这变化——周遭世界里的、他们自己的微观世界里的、以及他们自身之中的变化。我也说不清这奇怪的变化是何时、如何发生的——它想必是悄悄来的,蹑手蹑脚:有名望的人是叫人怕的。我自己就叫人生畏——不是叫我的学生或朋友怕,也不是叫那些亲自认识我的人怕,至少是叫那些只通过名声认识我、而自己并未感到有相当的名声庇护的人怕。

我是在将近十五年前那次“觉醒”之后,才意识到数学世界里(在其他科学圈子里同样如此,如果不是更甚)盛行的恐惧的。在此前的十五年里,我逐步地、不知不觉地进入了“大老板”的角色,进入了数学界《名人录》的世界。也同样是不知不觉,我成了这个角色的囚徒,它把我与所有人隔开,只除了少数几位“同侪”和少数几个学生(而且还……)——那些确实“有心气”的人。只有当我走出这个角色,环绕着它的恐惧至少有一部分才落了下来。那些在我面前多年沉默的舌头,松开了。

它们带给我的见证,不只是恐惧的见证,也是轻
% ---------------- p.143 ----------------
\origpage{24} 蔑的见证。尤其是身居其位的人对别人的轻蔑——一种唤起恐惧、也滋养恐惧的轻蔑。

我几乎不曾体验过恐惧,却实实在在体验过轻蔑——那是在一个人的存在与生活分量很轻的年月。我本来乐于忘掉轻蔑的时代,可它偏偏来向我问好了!也许它从未停止过,只是我满足于换了个世界(在我看来如此),满足于把目光移开,或者干脆:除了那些激动人心、没完没了的数学讨论之外,装作什么也看不见、什么也听不见?那些日子里,我终于肯去了解:轻蔑就在我周围到处盛行,在我选为自己的世界、我所认同、我所背书、并且曾娇养我的那个世界里盛行。

上面这些话也许给人一个印象:我被那些几乎一夜之间涌向我的见证搅得心神不宁。其实不然。这些见证是在一个仍然浮浅的层面上被接收的。它们只是加在我刚得知的另一些事实上,或加在我早已知道、却一直避免去注意的那些事实上。今天,我会这样表述当时学到的教训:“科学家”们,从最显赫的到最无名的,都是些和所有别人一模一样的人!我曾乐于想象“我们”是更好的一类,我们身上多点什么——我足足花了一两年,才摆脱这个顽固透顶的幻觉!

\section{(11) 与克劳德·谢瓦莱相遇——或自由与善意}

在帮助我做到这一点的朋友中,只有一位属于那个我刚离开、不打算再回去的圈子\fnoteP{4}{(6)}{\emph{我在“生存与生活”的朋友们}\par 在这些朋友中,我大概还该算上皮埃尔·萨米埃尔(Pierre Samuel)——我先前主要是在布尔巴基里认识他的,与谢瓦莱一样;而他(和谢瓦莱一样)在“生存与生活”小组里起了重要作用。我不觉得萨米埃尔有那么迷恋那种关于科学家优越性的幻觉。我感到,他带来最多的,是他在共同工作、讨论和待人接物中投入的常识和那份笑呵呵的好脾气,还有他优雅地在一个倾向于激进分析与激进抉择的团体里扮演“讨厌的改良派”的角色。在我退出“生存与生活”之后,他又在那里待了一阵,充当同名公报的主编;当他感到自己的存在对这个团体已不再有用时,便体面地离开了(去了“地球之友”)。\par 萨米埃尔和我出身于同一个狭小的圈子,这并不妨碍他成为那些沸腾年月里的朋友之一——我相信我从他们那里学到了些东西(尽管我是个糟糕的学生……)。他的为人方式,和谢瓦莱的为人方式一样——尽管二人毫不相像——对我那些“功绩主义”的倾向来说,是比最锋利的分析更好的解药!\par 现在我觉得,对那个时期所有让我学到点东西的朋友,我学到的主要不是通过解释、讨论之类,而是通过他们各自的为人方式、他们那与我不同的敏感——而“某种东西”最终就这样传达了过来……说到这里,我特别记得的,除了谢瓦莱和萨米埃尔,还有德尼·盖吉(Denis Guedj,他对“生存与生活”小组有很大的影响力)、达尼埃尔·西博尼(Daniel Sibony,他一直与这个小组保持距离,同时从角落里用一半鄙夷、一半打趣的眼光注视着自己的演变)、戈登·爱德华兹(Gordon Edwards,他是1970年6月在蒙特利尔那场“运动”诞生的共同行动者,并且多年间付出了惊人的精力,维持着《生存与生活》公报的英文“美国版”)、让·德洛尔(Jean Delord,一位与我年纪相仿的物理学家,精细而热情,他喜欢我,也喜欢“生存”这个小圈子)、弗雷德·斯内尔(Fred Snell,另一位在美国布法罗的物理学家, 1972年我在那里逗留了几个月,曾在他乡间的家中做客)。\par 在所有这些朋友中,五位是数学家,两位是物理学家,而全都是科学家——这似乎表明,那些年里与我最接近的圈子仍然是科学家的圈子,尤其是数学家的圈子。} (6)。他就是克劳德·谢瓦莱。他不发议论,对我的议论也不感兴趣,然而我想我可以说:我从他那里学到的东西,比我刚才说的那些更重要、也更隐蔽。

在我相当经常地与他来往的那段日子(“生存”小组的日子,他带着打了折扣的信念加入了这个小组),他常让我不知所措。我说不清是怎么回事,但我感到他握有某种我抓不到的
% ---------------- p.144 ----------------
\origpage{25} 认识——对某些想必极其简单、极其本质之物的理解;这些东西当然可以用简单的话说出来,但理解并不会因此就从一个人“传”到另一个人那里。我现在明白,他和我之间有一种成熟程度上的差别,这使我常常在他面前觉得自己站得不稳,像在进行一场聋子的对话——而这并不是因为彼此缺少好感或敬意。

虽然他没有用这样的词说过(就我记忆所及),但在他看来想必很清楚:我当时——或独自、或凭“生存”小组(后来成了“生存与生活”)内部共同的思考与活动的逻辑——达到的那些“质疑”(关于“科学家的社会角色”、关于科学等等),骨子里仍然是浮浅的。它们当然涉及我生活的世界,甚至涉及我在其中扮演的角色——但并没有真正深刻地牵动我。在那沸腾的年月里,我看待自己的眼光一丝未变。我并不是在那时才开始认识自己的。只是六年之后,我才生平第一次摆脱一个顽固的幻觉——不是关于别人或周围世界的幻觉,而是关于我自己的。那是另一次觉醒,分量比预备了它的第一次更大。它是一连串相继“觉醒”的“瀑布”中最初几次之一;我希望,这瀑布在我余年里还会继续流下去。

我不记得谢瓦莱曾在什么场合提到过认识自己,或者更确切地说,“发现自己”。不过回过头看,很清楚:他准是很早以前就已开始认识自己了。他有时也谈起自己,只在谈这谈那时说上几个字,那份朴素叫人不知所措。他是我没听见说过一句套话的两三个人之一。他话不多,而他说的话所表达的,不是他采纳并据为己有的观念,而是对事物的亲身感知与理解。所以他常叫我不知所措——早在我们还在布尔巴基小组里碰面的时候就是如此。他所说的,往往冲撞了我所珍视的、并因此被我当作“真”的那些看法。他身上有一种我所缺乏的内在自主,我在“生存与生活”的日子里开始隐约察觉它。这种自主不属于智性、不属于话语的层面。它不是一件可以像观念、观点之类那样“采纳”的东西……幸好,我从未想过要把在别人身上察觉到的这份自主“据为己有”。我得找到我自己的自主。这同时意味着:我要学会(或重新学会)做我自己。但在那些年里,我对自己的不成熟、对自己缺乏内在自主,毫无觉察。我若最终发现了这一点,那么与谢瓦莱的相遇,必定是那些在我身上默默发酵的酵母之一——那时我正投身于一些大计划。播下这酵母的,不是议论,也不是言语。要播下它,只需我在路上偶然遇见的某个人不讲议论,只满足于做他自己。

我觉得,在七十年代初,当我们因《生存与生活》公报的出版而定期见面时,谢瓦莱曾不加坚持地试图向我传达一个消息,而我当时太笨手笨脚,抓不住它,或太埋头于我的活动任务。我隐约意识到,他有关于自由——关于内在自由——的东西要教我。我那时倾向于靠一些宏大的道德原则运转,而且从《生存》最初几期起就开始吹那把小号,觉得理所当然;而他对说教式的话语有一种特别的反感。我想,在《生存》创办之初,这是我最摸不透他的地方。在他看来,这种话语只是一种强制的尝试,叠加在窒息人的无数别的外在强制之上。当然,一个人尽可以花一辈子去讨论这样
% ---------------- p.145 ----------------
一种看法,讨论正反两面\origpage{26}。它完全冲撞了我的看法——那看法(可想而知)是由最高尚、最普遍的感情驱动的。我很难过:我无法理解,谢瓦莱——我最敬重的人,我多少把他当作战友的人——竟然会带着一种狡黠的乐趣,不去分享这些感情!

\origpage{27} 我不明白:真相、事物的实际,既不是善良感情的问题,也不是观点或偏好的问题。谢瓦莱看见了一件东西——那东西再简单、再真实不过——而我没有看见。并不是他在什么地方读到过它;看见一件东西,和读到关于它的东西,二者毫无共同之处。人勉强可以用手(盲文)或耳朵(若有人读给你听)去“读”一段文字,但要看见那件东西本身,只能用自己的一双眼睛。我不认为谢瓦莱的眼睛比我的好。但他用它们,我不用。我被我的善良感情之类的东西占得太满,没有闲工夫去看我的善良感情和原则对我自己、对别人——首先是对我自己的孩子们——起了什么效果。

他一定看得很清楚:我常常不用自己的眼睛,甚至一点也不想用。奇怪的是,他从未对我透露过这一点。还是他说了,只是我没听见?还是他忍住了,认为那是白费力气?又或者他压根儿没想到要说——说到底,我用不用自己的眼睛,是我的事,不是他的事!

\section{(12) 功绩与轻蔑}

我想借我自己有限的经验细看一下:轻蔑是何时、如何进入数学家世界的,尤其是在那个由同事、朋友和学生构成的“微观世界”里——它已成了我的第二祖国。同时我也想看看,在这场转变中,我这一份占了多少。

我想我可以毫无保留地说: 1948–49年间,在我前面谈过的那个数学家圈子里(对我来说,它的中心是最初的布尔巴基小组),我没有遇到一丝一毫的轻蔑,甚至没有一丝一毫的鄙夷或居高临下——无论对我自己,还是对任何一个来此学习数学家行当的年轻人,无论法国人还是外国人。在那里起旗舰作用的人物——凭地位或声望,如勒雷、嘉当、韦伊——我并不怕他们,我想我的同学里也没有人怕他们。勒雷和嘉当很有“体面先生”的派头,另当别论;至于其余那些随随便便闯进来、像哥们儿一样跟嘉当“你你我我”、显然又“在行”的家伙,我甚至花了相当一段时间才明白:他们每一个都和嘉当本人一样是大学教授,绝不像我那样挣一天吃一天,而是领着在我看来属于天文数字的薪水,而且还是个有国际影响、国际听众的数学家。

听从韦伊的建议,我在南锡度过了随后的三年。那里当时算是布尔巴基的一个总部:德尔萨特、迪厄多内、施瓦茨、戈德芒(稍后还有塞尔)都在大学里任教。跟我一起在那里的只有四五个年轻人(我记得有利翁斯、马尔格朗热、布吕阿、贝热,若有记错请原谅),所以远不像在巴黎那样“淹在人堆里”。气氛也因此更亲切,大家彼此都认识,我想彼此都以“你”相称。可我翻检记忆时,我亲眼看见一个数学家毫不掩饰地轻蔑对待一个学生的头一遭、也是唯一一遭,恰恰就在那里。那可怜人是当天从另一座城市赶来的,要跟他的老板一起工作。(他当时应该在准备博士论文,后来还体面地通过了,而且此后也
% ---------------- p.146 ----------------
有了相当的名声,我想。)\origpage{28} 那场面让我相当震惊。要是有人敢用这种腔调跟我说哪怕一秒钟的话,我立马就把他摔门轰出去!可这一回,我熟的只是那位“老板”,甚至跟他称兄道弟,而不是那个学生——我只是见过他。我这位年长的同行,除了广博的学养(不只是数学的)和锋利的头脑,还有一种斩钉截铁式的权威感,在当时(而且在之后相当长一段时间里,直到七十年代初)都令我敬畏。他对我有某种影响力。我不记得自己有没有就他的态度问过他什么;我只记得我从那场面里得出的结论:这个可怜学生一定是差劲透了,才该受这样的对待——诸如此类。我当时没有对自己说:这学生若是真的差劲,那该劝他去做别的,该停止跟他合作,但无论如何不该轻蔑地对待他。我认同了像这位显赫前辈那样的“数学强者”,而以那些“废物”为代价——似乎轻蔑他们是正当的。于是我走上了那条早已画好的路:与轻蔑默契,因为这对我方便,它还突出了这一点——我,被接纳入有功者、“数学强者”的兄弟会了!\fnoteP{5}{(7)}{上面这一段,是整篇导论里在我的原稿上涂改得最厉害、增补最多的一段。对那件事的描述、用词的选择,最初都是顶着来的,逆着潮流——显然有一股力量在推着人快点把那件事带过去,好像只为对良心有个交代,好“转到正经事上去”。这些都是一份抵抗的熟悉征兆——这里抵抗的是对这段插曲的澄清,以及对它作为某种内心态度之揭示者的意义的澄清。情形与这篇导论开头(第2段)所描述的十分相似:在数学工作中发现一处矛盾及其意义的那个“关键”时刻——那时起作用的是精神的惰性,是它不肯与一种错误或不充分的理解(但其中根本没有我们这个人牵涉在内)分离的厌恶。后者却是主动的,必要时还能别出心裁,做到无水也能把鱼淹死;而我刚才说的惰性只是一种被动的力量。在现在这个情形里,比在数学工作的情形里更甚:刚刚以其全部的简单、全部的自明显现出来的发现,紧接着就是一阵卸下重负的轻松感,一种解放感。这不止是一种感觉——不如说是对刚刚发生之事的一种敏锐而感激的知觉,那是一次解放。} (7)

当然,我不会比任何人更清楚地对自己说:那些试着搞数学却搞不成的人,活该被轻蔑!如果我在那个时候或任何时候听见有人说这种话,我会漂亮地驳他一顿,真心为如此惊人的精神无知而难过。事实是:我泡在暧昧里,同时在两张互不相通的牌桌上玩:一边是漂亮的原则与感情,另一边是:可怜的家伙,得差劲成什么样才落得这般对待(言下之意:这种事肯定不会落到我头上,那当然!)。

最后我觉得,我讲的那桩事件,尤其是我在其中扮演的那个(看来无伤大雅的)角色,其实是我身上一种暧昧的典型;这暧昧在此后二十年里跟了我整个数学家的生涯,直到1970年“觉醒”之后才消散\fnoteP{6}{(8)}{正如下文将变得清楚的,这种暧昧丝毫没有“在1970年觉醒之后消散”。这里有一个“自我”典型的战略撤退:它把“觉醒之前”那段时期丢给盈亏自负,而那段时期随即成了划分出一个无可指摘的“之后”的分界线!} (8)——而直到今天,写下这些字句的时候,我才清楚地察觉它。也真可惜,我当时竟没有察觉。也许对我来说时机尚未成熟。总之,那些当时传到我耳中的、关于轻蔑之统治的见证——对那统治,我选择了闭上眼睛——并没有把我个人牵涉进去,也没有把我那亲爱的微观世界里离我最近的那部分同行和朋友中的任何一位牵涉进去\fnoteP{7}{(9)}{这并不完全准确:在我最亲近的同行中至少有一个例外,后文会看到。这里有一种记忆特有的“懒”,它往往倾向于把那些与由来已久的、根深蒂固的看法“不合”的事实“随手放过”。} (9)。它们更像是这个调子:唉!竟然不得不%
% 原句在此页末尾中断, 下文接 part1-ch06b.tex (p.147)
